\documentclass[11pt,a4paper]{article}
% Shared preamble for RuPhO English problem statements (match T1 2026 style).
% Use: \input{latex/rupho_en_preamble.tex} after \documentclass.
\usepackage[margin=1in]{geometry}
\usepackage{amsmath,amssymb}
\usepackage{graphicx}
\usepackage{float}
\usepackage{enumitem}
\usepackage{microtype}
\usepackage{caption}
\usepackage{tabularx}
\usepackage{hyperref}

\captionsetup{font=small,labelfont=bf,skip=6pt}

\setlist[itemize]{leftmargin=1.2cm}
\setlist[enumerate]{leftmargin=1.2cm}

% One outer box per subproblem: [id | statement | points] (no inner rules)
\newcommand{\problembox}[3]{%
  \par\addvspace{0.42em}%
  \noindent
  \setlength{\fboxsep}{7.5pt}%
  \setlength{\fboxrule}{0.95pt}%
  \fbox{%
    \begin{minipage}{\dimexpr\linewidth-2\fboxsep-2\fboxrule\relax}
    \begin{tabularx}{\linewidth}{@{}>{\raggedright\arraybackslash\bfseries}p{2.1em} @{\hspace{0.9em}} >{\raggedright\arraybackslash}X @{\hspace{0.9em}} >{\raggedleft\arraybackslash\bfseries}p{2.4em}@{}}
    #1 & #2 & #3
    \end{tabularx}
    \end{minipage}%
  }%
  \par\addvspace{0.32em}%
}


\title{\textbf{Half empty bottle}}
\author{\normalsize Y21 (2021) --- English translation}
\date{}
\begin{document}
\maketitle

When we push a half-empty bottle of water along a plane, vibrations of the liquid inside it cause a sharp jerking movement of the system. In this task we will study this process and build a theory to explain this phenomenon. To get you started, let's look at the results of an experiment in which we push a bottle of water and see how its position depends on time. In all experiments during the task, the same bottle will be used, so its parameters data in one part can be used in another.

\begin{figure}[H]
    \centering
    \includegraphics[width=0.72\textwidth]{PE1_EN_assets/fig1.jpg}
    \caption{Figure 1.}
\end{figure}

\section*{Part A. Kinematics of water in a bottle}

First of all, let's look at how water flops around inside a cylindrical vessel. For experimental study, let's fix a bottle of radius $R=5.0~\text{cm}$ on a horizontal axis and, by means of small translational oscillations of the axis, create oscillations in the water level, and then wait until the system moves into some fairly stable stationary state of liquid oscillations inside the vessel. To measure the speed of liquid flow, we fill it with polystyrene balls (density $1.06~\text{г}/\text{cm}^3$, diameter $0.7~\text{mm}$) and photograph the stationary state of the system with a camera with a shutter speed (shutter speed is the time interval during which light hits an area of ​​the photosensitive material or photosensitive matrix) $T_\text{e}=1/10~\text{с}$, this time is less than the half-cycle of water oscillations in the vessel.

\begin{figure}[H]
    \centering
    \includegraphics[width=0.72\textwidth]{PE1_EN_assets/fig2.jpg}
    \caption{Figure 2.}
\end{figure}

The angular momentum of the water relative to the axis of the bottle is related to the angular velocity of rotation of the water plane $\dot{\theta}$ and the characteristics of the system through a certain dimensionless coefficient $\alpha$: \textbackslash\{\}[ L_\textbackslash\{\}text\{w\} = \textbackslash\{\}alpha \textbackslash\{\}cdot \textbackslash\{\}pi \textbackslash\{\}rho R^4 L \textbackslash\{\}dot\{\textbackslash\{\}theta\}, \textbackslash\{\}] where $\rho$ is the density of water, $R$ is the radius of the bottle, $L$ is its length.

\problembox{A1}{What is the theoretical value of $\alpha_\text{solid}$ if we replace water with a solid body of constant density coinciding with the density of water?}{0.20}

\problembox{A2}{Determine the coefficient $\alpha$ in the experiment performed. Specify the points at which you measure the speed with a pair of numbers in polar coordinates centered on the axis of the bottle $r, \varphi$.}{4.50}

\problembox{A3}{Construct a graph characterizing your measurements in coordinates $r/R$ vs $\varphi/\pi$.}{0.50}

\problembox{A4}{Compare the resulting $\alpha$ with the theoretical estimate: \textbackslash\{\}[ \textbackslash\{\}alpha_\textbackslash\{\}text\{theory\} = \textbackslash\{\}frac\{4\}\{\textbackslash\{\}pi^2\} - \textbackslash\{\}frac\{1\}\{4\}, \textbackslash\{\}] obtained in the 19th century by Lord Rayleigh.}{0.20}

Now we will measure how the angle $\theta$ depends on time during stationary oscillations. Next, we apply the Fourier transform to the obtained dependence $\theta(t)$, i.e. we obtain the spectral characteristics of free vibrations of the liquid inside the bottle. Designate the distance from the center of mass of the half-cylinder to its axis as $a$.

\begin{figure}[H]
    \centering
    \includegraphics[width=0.72\textwidth]{PE1_EN_assets/fig3.jpg}
    \caption{Figure 3.}
\end{figure}

\problembox{A5}{Based on the data presented in the figure, estimate the system’s own oscillation frequency $\omega_0$ and its quality factor $1/\zeta$.}{0.30}

\problembox{A6}{Estimate $\alpha$ from the spectral response of the vibrations.}{0.70}

\section*{Part B. System Dynamics}

In the last part, we came to the conclusion that the vibrations of water in bottles can be approximated with acceptable accuracy by the vibrations of a semi-cylindrical solid body with a moment of inertia $I_0 = \alpha \cdot \pi \rho R^4 L$ and a mass $m=\dfrac{\pi}{2} \rho R^2 L$ around its axis. Thus, a bottle of water can be represented as a system of rings of total mass $M$ and a half-cylinder coaxial with them. Moreover, the rings of mass $M$ roll without slipping along the plane, and the half-cylinder is suspended by the axis of the rings and does not interact with the plane.

\begin{figure}[H]
    \centering
    \includegraphics[width=0.72\textwidth]{PE1_EN_assets/fig4.jpg}
    \caption{Figure 4.}
\end{figure}

The entire system can move along the horizontal axis $x$ perpendicular to the axis of rotation of the system. The parameters of our system will be a pair of variables $x$ and $\theta$, so write all equations in your answers in terms of these variables and their time derivatives, for example $\ddot{x}$.

\problembox{B1}{Write down Newton's II law for the entire system and the equation of rotational motion for the rings.}{0.30}

\problembox{B2}{From the equations of the previous paragraph, obtain an equation that contains only constants, parameters of our system and their derivatives.}{0.20}

\problembox{B3}{Write down the equation for the dynamics of rotational motion for a half-cylinder as a second-order differential equation.}{0.30}

\problembox{B4}{Find the angular frequency of oscillations of the system $\Omega$. Write down the solutions to the resulting system of two equations, which describes how the bottle rolls jerkily along the plane.}{0.50}

To confirm this theory, we will conduct the following experiment. Let us progressively accelerate the system to speed $v_0$ and collide it with the wall. In our case, the mass of water is $m=2.780~\text{кг}$, $M=0.910~\text{кг}$, $L=70~\text{cm}$. Since the mass of water is several times greater than the mass of the bottle, we will assume that $\eta$ of kinetic energy before the impact turns into vibrational energy, and the rest is dissipated. Below is a diagram of how this happens.

\problembox{B5}{Let the amplitude of the bottle speed oscillation $\gamma v_0$ after rebounding from the wall. Express $\gamma$ in terms of known quantities and $\eta$.}{0.30}

The following are the measurement results for $\dot{x}/v_0$ vs $t$.

\begin{figure}[H]
    \centering
    \includegraphics[width=0.72\textwidth]{PE1_EN_assets/fig5.jpg}
    \caption{Figure 5.}
\end{figure}

$t,~\text{с}$ $\dot{x}/v_0$ $t,~\text{с}$ $\dot{x}/v_0$ $t,~\text{с}$ $\dot{x}/v_0$ -0.75 -0.96 0.55 0.42 1.42 0.41 -0.66 -0.98 0.58 0.26 1.47 0.23 -0.53 -0.96 0.61 0.10 1.51 0.08 -0.41 -0.96 0.63 0.03 1.55 -0.09 -0.31 -0.96 0.64 0.04 1.60 0.00 -0.19 -0.97 0.70 0.20 1.65 0.29 -0.18 -0.51 0.73 0.34 1.67 0.42 -0.14 0.00 0.76 0.47 1.71 0.45 -0.08 0.01 0.78 0.52 1.74 0.38 -0.02 0.01 0.83 0.50 1.79 0.14 0.04 0.07 0.90 0.25 1.84 0.00 0.07 0.19 0.92 0.04 1.86 -0.03 0.08 0.31 0.95 -0.07 1.92 0.07 0.13 0.44 0.96 -0.05 1.93 0.22 0.17 0.57 1.02 0.13 1.96 0.37 0.23 0.47 1.05 0.26 2.03 0.52 0.26 0.34 1.07 0.41 2.06 0.43 0.28 0.18 1.10 0.41 2.08 0.30 0.31 0.08 1.16 0.22 2.15 0.06 0.33 0.05 1.21 -0.04 2.18 0.02 0.36 0.13 1.24 -0.12 2.18 0.21 0.39 0.33 1.31 0.09 2.18 0.03 0.43 0.50 1.34 0.26 2.27 0.36 0.49 0.59 1.37 0.37 2.31 0.51 0.50 0.57 1.40 0.41 — —

\problembox{B6}{Plot $\dot{x}/v_0$ vs $t$.}{1.50}

\problembox{B7}{Find the experimental value of $\Omega$, compare it with the theoretical one.}{0.20}

\problembox{B8}{Find the value of $\eta$.}{0.30}

\vspace{1em}
\noindent\textit{Source:} \url{https://pho.rs/p/238}

\end{document}